\section{Gemini 2.5 Flash: capacidad de visión multimodal nativa}

\subsection{Detección, segmentación y reconocimiento conversacionales}

Gemini 2.5 Flash representa una ruta más «de extremo a extremo»: no necesita modelos especializados externos como SAM3, pues el modelo mismo cuenta con la capacidad de detección, segmentación y reconocimiento.

\begin{importantbox}{Capacidad de visión nativa de Gemini 2.5 Flash}
Hasta diciembre de 2025, el único VLM que admite una salida de segmentación Mask nativa es Gemini 2.5 Flash. Sus características clave:
\begin{itemize}
    \item Interacción conversacional: la segmentación se activa mediante una Query en lenguaje natural
    \item Formato de salida: List of Dict, donde cada Dict contiene \texttt{box\_2d} (cuadro delimitador), \texttt{mask} (imagen de la máscara de segmentación), \texttt{label} (etiqueta descriptiva)
    \item Las coordenadas de BBox están normalizadas dentro de un rango de 1000×1000 y deben convertirse al superponerlas sobre la imagen original
\end{itemize}
\end{importantbox}

\subsection{Forma de uso}

Configuración clave para invocar a Gemini 2.5 Flash y realizar la segmentación:

\begin{itemize}
    \item Usar el modelo \texttt{gemini-2.5-flash}
    \item Establecer \texttt{temperature=0} (para garantizar una salida determinista)
    \item Pasar un prompt de imagen y texto, usando la plantilla de prompt de segmentación recomendada oficialmente
\end{itemize}

\subsection{Evaluación del desempeño de la segmentación}

Tomando como ejemplo una imagen de prueba con varios niños corriendo:
\begin{itemize}
    \item De los seis niños, la detección y segmentación de cinco de ellos es muy precisa
    \item El cuadro del tercer niño presenta una ligera desviación
    \item La precisión general es muy alta; la calidad de la Mask corresponde a una segmentación a nivel de píxel mediante una matriz 0/1
\end{itemize}

\begin{knowledgebox}{Diferencia entre Gemini 2.5 Flash Preview y Pro}
\begin{itemize}
    \item \textbf{Gemini 2.5 Flash}: admite la salida completa de Mask; es actualmente el único VLM que admite la segmentación conversacional nativa
    \item \textbf{Gemini 2.5 Flash Preview} (recién publicado): por ahora no admite la salida de Mask
    \item \textbf{Gemini 2.5 Flash 3 Pro}: es posible que solo genere puntos de contorno en lugar de una Mask completa
\end{itemize}
\end{knowledgebox}

\subsection{Resumen de la sección}

Gemini 2.5 Flash muestra una dirección importante en el desarrollo de los VLM: interiorizar de forma nativa las capacidades de detección, segmentación y reconocimiento en un modelo unificado, y lograr así una auténtica «comprensión visual conversacional». Aunque la combinación de SAM3 + VLM resulta más flexible, el esquema de extremo a extremo representa la tendencia del futuro.


